% !TeX root = ../papers/nc-p3.tex
\section{節点・尖点三次曲線上の順序付き三交換子の全逆像}
\label{sec:marked-singular-integral-commutators}
\begin{theorem}\label{thm:marked-singular-integral-commutators}
標数零の代数閉体上、順序付き跡零三行列の鉛筆 $C(v)$ の行列式が
特異な整三次曲線 $E$ を定めるとする。特異点で $C$ の階数が1なら
$C=C_Q$ を満たす順序付き跡零三行列は存在しない。階数2なら次が全逆像である。
\begin{enumerate}
\item 節点型では四つの孤立解、二つの孤立解、あるいは二つの孤立解と一つの
滑らかなアフィン二次曲面 $uw-v^2=1$ となる。一般には四解である。
\item 尖点型では二つの孤立解、空集合、あるいは上と同じ一つのアフィン二次曲面となる。
一般には二解である。
\end{enumerate}
各型の任意の局所自由な目標表示類で、掲げた三通りが全て起こる。
以下の線型代数と余因子方程式で、基底を固定した入力ごとに区別し、全ての解を与える。
\end{theorem}
\begin{proof}[整三次曲線上の補助表示の被覆]
定理\ref{thm:marked-elliptic-commutator-doubling}の
$W=\operatorname{im}\Gamma_Q$ と六次元線型解の証明は整な行列式にそのまま適用できる。
補助表示 $A$ の左右基底同値類を調べる。転置の余核層 $T$ は $E$ 上の階数一捩れなし層である。
実際、平面上の長さ一の自由分解により各閉点で深さ1となる。
分解から $\chi(T(n))=3n+3$、$h^0(T)=3$、$h^0(T(-1))=0$ である。
正規化 $\pi:\mathbf P^1\to E$ で局所自由でない場合は
$T=\pi_*\mathcal O_{\mathbf P^1}(2)$ しかない。特異点の完備局所環を $R$、正規化を
$\bar R$ とする。分数イデアルとしての局所加群を $I\bar R=\bar R$ に規格化すると、
導手 $\mathfrak c$ は $I$ に入る。$I/\mathfrak c$ は二次元の $\bar R/\mathfrak c$
（節点では $k\oplus k$、尖点では $k[\epsilon]/(\epsilon^2)$）をイデアルとして生成する
部分空間である。一次元なら単元で生成され $I$ は自由、二次元なら $I=\bar R$ である。
後者では正規化上の線束の次数がEuler標数から2になる。
局所自由の場合は正規化で次数3、節点の二値の貼合せは乗法群、
尖点の一次ジェットの貼合せは加法群で指定される。$h^0(T(-1))=0$ により恒等元を除く。

節点を $xyz=y^3+z^3$、正規化を $(r^3+1:r^2:r)$ とすると、右核点は
$(r^3+\lambda:r^2:r)$ であり $\lambda\in k^*\setminus\{1\}$ である。
非自由表示は $\lambda=0$ で共通因子を除いたものとなる。双一次関係の基底は
\[
 N_\lambda=\begin{pmatrix}0&-z&y\\-y&x&(\lambda-1)z\\
 -z&-(\lambda-1)y&\lambda x\end{pmatrix}.
\]
尖点では $xz^2=y^3$、正規化 $(1:r^2:r^3)$ を取り、右核点は
$(1+\lambda r:r^2:r^3)$、$\lambda\ne0$、非自由なら $(1:r:r^2)$ である。
$r$ の倍率による生成添字の基底変更で、非零の $\lambda$ を1にしてよい。二表示は順に
\[
 K=\begin{pmatrix}y&-x&-x\\z&-y&-x\\0&z&-y\end{pmatrix},\qquad
 K_0=\begin{pmatrix}y&0&-x\\z&-y&0\\0&z&-y\end{pmatrix}.
\]
実際各行は正規化上のグラフを消す。自由な三次表示では一次の双一次関係は自由分解により三本で、
その全空間が元の行空間を復元する。明示した各代表の行は独立で、行列式は $E$ の式の非零倍である。
以上は目標 $C$ 自身の左右基底同値類にも適用できる。
\end{proof}

\begin{proof}[負二倍、六座標と余因子の計算]
対称行列 $Z=\left(\begin{smallmatrix}P&A&B\\ A&T&D\\ B&D&U\end{smallmatrix}\right)$ を取る。
六空間 $\Gamma_Q\subset W$ のうち、三行列の次の六成分を指定して $Q(Z)$ を一意に定める。
成分の添字は零始まりである。節点では
\[
 ((Q_1)_{02},(Q_2)_{01},(Q_2)_{02},-\lambda(Q_1)_{21},(Q_2)_{12},(Q_1)_{12})
 =
 (-\lambda((\lambda-1)A+U),-(\lambda-1)B+\lambda T,
 (\lambda-1)^2P+\lambda D,-\lambda A,-\lambda B,(\lambda-1)P).
\]
尖点 $K$ では第四成分を $-(Q_1)_{21}$、第六成分を $(Q_2)_{10}$ とし、
右辺を $(U,-A-B-T,-B-D,A,B,P)$ とする。
これらを代入して六空間の線型方程式を解くと各成分は $Z$ と係数 $\lambda$ の多項式
（定数分母3を許す）である。六成分指定の行列式は符号を除き節点で
$\lambda^5(\lambda-1)$、尖点で1となるので、六次元性により全解が得られる。
節点の参照行列と尖点の参照行列を順に
\[
 M_\lambda=N_{1/\lambda^2}\operatorname{diag}(1,\lambda,\lambda^2),\qquad
 M_K=\begin{pmatrix}-y&x+y&-x+y\\ y+z&-x+y&0\\-y&x+y+z&-x+z\end{pmatrix}
\]
とする。尖点の右核は $((r-1)^2:r^2-r^3:r^3)$、従って $1-2r,r^2,r^3$ の別基底である。
以上の右核の貼合せは節点で $\lambda\mapsto\lambda^{-2}$、尖点で
$\lambda\mapsto-2\lambda$ となる。後者は $r$ の倍率を元に戻した式である。

$\lambda\ne\pm1$ の節点および自由な尖点で、線型同型
$\Psi:\mathrm{Sym}_3\to\{L:\operatorname{tr}(LM(v))=0\ \forall v\}$ は
\[
                   C_{Q(Z)}(v)=\Psi(\operatorname{adj}Z)M(v)
\]
で計算できる。検算の手順を示す。$C_{Q(Z)}(0,1,0)M(0,1,0)^{-1}$ の
$TU,PU,PT,BD,AD,AB$ の係数を順に六列とする。
$\operatorname{adj}Z$ の対角三成分と上三角非対角三成分 $01,02,12$ をこの六列に掛けると、
$Z$ の全二次係数が元の左倍率と一致する。他の二スライスも同じ左倍率と $M$ の積である。
六列の成分を行優先に並べ、節点では行 $0,1,2,3,4,5$、尖点では行 $0,1,2,3,4,7$ を取ると
小行列式はそれぞれ $\lambda^2(\lambda-1)(\lambda+1)$ と $-3$ である。
節点の六列の分母は $\lambda$ の冪だけなので恒等式を $\lambda=-1$ に特殊化できる。
そこで $M_{-1}$ は行列式が恒等零ゆえ補助解は一つも整な目標を与えない。
非自由表示 $N_0,K_0$ の双方では六空間の任意の二基底要素 $q,q'$ について
$[q_1,q'_2]+[q'_1,q_2]=0$ となる。この場合 $C_Q(1,0,0)=0$ で行列式は二変数のみになる。
これらの全係数の厳密検算は \path{computations/verify_singular_integral_commutators.py} にある。
従って非自由な目標に至る補助表示はなく、自由な目標の補助類は上の負二倍の逆像で尽きる。
\end{proof}

\begin{proof}[指定目標と全例外]
自由な目標を各候補について $C=L M R$ と書く。定理\ref{thm:marked-fixed-commutator-adjugate}と同じく
$S=\Psi^{-1}(RL)$ とすれば、解は $\operatorname{adj}Z=S$ の各逆像につき
$R^{-1}Q(Z)R$ である。ここでも右核像の全点を固定する射影変換は恒等であり、他の共役はない。
$S$ の階数3,2,1に対応する解は二点、空、一つのアフィン二次曲面である。

節点の二選択の関係を明示する。$M_{-\lambda}=M_\lambda\operatorname{diag}(1,-1,1)$ ゆえ、
一方の $S$ の座標を $(P,T,U,A,B,D)$ とすれば、他方の $S'$ の座標は
$(P-4D,T,U,2U-A,2T-B,-D)$ となる。先の六列への代入で
\[
          4\det\Psi_\lambda(S)=(\lambda+1)^2\det S-(\lambda-1)^2\det S'
\]
を得る。左辺は零でないので少なくとも一方の $S$ は可逆である。
両行列式の非消失は稠密開条件である。$S=E_{11}$（中央の対角成分、零始まり）
では $\det S'=-4$ なので、階数1の例を得る。階数2の軌跡は階数1を閉包に含むので
$\det S'\ne0$ の階数2の例もある。尖点では同じ $S=E_{11}$ で $\det\Psi(S)=1$、
やはり階数2および3に動かせる。これで固定した各自由表示類における全ての可能性が従う。
\end{proof}
